\SourceLeaf{77}%
\begin{center}
{\large THEORY OF BIQUADRATIC RESIDUES.\par}
\bigskip FIRST PAPER.
\end{center}
\BookRule{28mm}

\Article{1}
The theory of quadratic residues reduces to a few fundamental theorems that rank among the most beautiful jewels of higher arithmetic. As is well known, these were first easily discovered through inductive observation and then proved in many different ways, leaving nothing further to be desired.

The theory of cubic and biquadratic residues, however, requires much deeper investigation. When we began to explore it in 1805, besides the results that lie, as it were, at the threshold, some special theorems presented themselves, remarkable both for their simplicity and for the difficulty of their proofs. But we soon found that the principles of arithmetic used up to then were wholly insufficient to establish the general theory. Rather, the theory necessarily requires the field of higher arithmetic to be enlarged, as it were, infinitely many times over. What this means will become entirely clear as these investigations continue. As soon as we entered this new field, inductive observation immediately opened a path to knowledge of the simplest theorems encompassing the entire theory. Their proofs, however, lay so deeply hidden that only after many unsuccessful attempts could they finally be brought to light.

Now that we are preparing to publish these studies, we shall begin with the theory of biquadratic residues. In this first paper\space\SourceLeaf{78}%
 we shall present the investigations that could already be completed on this side of the enlarged field of arithmetic. They pave the way to that field, as it were, and also contribute some new developments to the theory of the division of the circle.

\Article{2}
We introduced the notion of a biquadratic residue in Article 115 of the \emph{Disquisitiones Arithmeticae}: an integer $a$, positive or negative, is called a biquadratic residue of the integer $p$ if $a$ can be congruent to a fourth power modulo $p$, and a biquadratic nonresidue if no such congruence exists. Throughout the following investigations, unless we expressly state otherwise, we shall assume that the modulus $p$ is an odd positive prime and that $a$ is not divisible by $p$, since all the remaining cases can very easily be reduced to this one.

\Article{3}
Clearly, every biquadratic residue of $p$ is also a quadratic residue of the same number, and hence every quadratic nonresidue is also a biquadratic nonresidue. The converse also holds whenever $p$ is a prime of the form $4n+3$. For if, in this case, $a$ is a quadratic residue of $p$, let us put $a\equiv bb\; (\mathrm{mod.}\,p)$. Here $b$ is either a quadratic residue of $p$ or a nonresidue. In the first case, let us put $b\equiv cc$, giving $a\equiv c^4$; that is, $a$ is a biquadratic residue of $p$. In the second case, $-b$ is a quadratic residue of $p$ (since $-1$ is a nonresidue of every prime of the form $4n+3$). Putting $-b\equiv cc$ again gives $a\equiv c^4$, and $a$ is a biquadratic residue of $p$. It is also easy to see that in this case the congruence $x^4\equiv a\; (\mathrm{mod.}\,p)$ has no solutions besides $x\equiv c$ and $x\equiv-c$. Since these evident propositions exhaust the entire theory of biquadratic residues for prime moduli of the form $4n+3$, we shall exclude these moduli altogether from our investigation, restricting it to prime moduli of the form $4n+1$.

\Article{4}
Thus, when $p$ is a prime of the form $4n+1$, the converse of the proposition in the preceding article does not hold: there can be quadratic residues that are not also biquadratic residues. This happens whenever a quadratic residue is congruent to the square of a quadratic nonresidue. For put $a\equiv bb$, with $b$ a quadratic \SourceLeaf{79}%
nonresidue of $p$. If the congruence $x^4\equiv a$ could be satisfied by a value $x\equiv c$, we would have $c^4\equiv bb$, so that the product $(cc-b)(cc+b)$ would be divisible by $p$. Then $p$ would have to divide either the factor $cc-b$ or the other factor $cc+b$; that is, either $+b$ or $-b$ would be a quadratic residue of $p$, and therefore both would be (since $-1$ is a quadratic residue), contrary to the hypothesis.

All integers not divisible by $p$ could therefore be divided into three classes: the first would contain the biquadratic residues, the second the biquadratic nonresidues that are also quadratic residues, and the third the quadratic nonresidues. Clearly it suffices to classify only the numbers $1,2,3\dotsfive p-1$. Half would belong to the third class, while the other half would be divided between the first and second classes.

\newpage\Article{5}
It will be better, however, to establish four classes as follows.

Let $A$ be the collection of all biquadratic residues of $p$ between $1$ and $p-1$, inclusive, and let $e$ be any chosen quadratic nonresidue of $p$. Further, let $B$ be the collection of the least positive residues of the products $eA$ modulo $p$, and let $C,D$ respectively be the collections of the least positive residues of the products $eeA$, $e^3A$ modulo $p$. It is then easy to see that the individual numbers in $B$ are distinct, as are those in $C$ and those in $D$; zero cannot occur among any of these numbers. Also, all numbers in $A$ and $C$ are quadratic residues of $p$, while all those in $B$ and $D$ are quadratic nonresidues. Thus neither of the collections $A,C$ can have a number in common with $B$ or $D$. But neither can $A$ have a number in common with $C$, nor $B$ with $D$. For suppose

I. that some number in $A$, say $a$, also occurs in $C$, where it arises from the congruent product $eea'$, with $a'$ a number in $A$. Put $a\equiv\alpha^4$, $a'\equiv\alpha'^4$, and choose an integer $\Theta$ such that $\Theta\alpha'\equiv1$. We then have $ee\alpha'^4\equiv\alpha^4$, and multiplying by $\Theta^4$ gives

\[
ee\equiv\alpha^4\Theta^4
\]

That is, $ee$ is a biquadratic residue, and hence $e$ a quadratic residue, contrary to the hypothesis.

\SourceLeaf{80}%
II. Similarly, suppose that a number belongs to both collections $B,D$, arising from the products $ea$, $e^3a'$, with $a,a'$ numbers in $A$. The congruence $ea\equiv e^3a'$ would imply $a\equiv eea'$. We would thus have a number arising from the product $eea'$ that belonged to both $C$ and $A$, which we have just proved impossible.

It is also easy to prove that \emph{all} quadratic residues of $p$ between $1$ and $p-1$, inclusive, must occur either in $A$ or in $C$, and all quadratic nonresidues of $p$ between those limits must occur either in $B$ or in $D$. For

I. Every such quadratic residue that is also a biquadratic residue belongs to $A$ by hypothesis.

II. Write a quadratic residue $h$ (less than $p$) that is also a biquadratic nonresidue as $\equiv gg$, where $g$ is a quadratic nonresidue. Choose an integer $\gamma$ such that $e\gamma\equiv g$. Then $\gamma$ is a quadratic residue of $p$, which we write as $\equiv kk$. Hence

\[
h\equiv gg\equiv ee\gamma\gamma\equiv eek^4
\]

Since the least residue of $k^4$ belongs to $A$, the number $h$, which arises by multiplying it by $ee$, must belong to $C$.

III. Let $h$ denote a quadratic nonresidue of $p$ between $1$ and $p-1$. Find an integer $g$ between the same limits such that $eg\equiv h$. Then $g$ is a quadratic residue and so belongs to either $A$ or $C$. In the first case, $h$ evidently belongs to $B$, and in the second to $D$.

All this shows that the numbers $1,2,3\dotsfive p-1$ are distributed among the four lists $A,B,C,D$ so that each occurs in one of them; each list must therefore contain $\tfrac14(p-1)$ numbers. In this classification, the numbers belonging to $A$ and $C$ are intrinsic to those classes. The distinction between $B$ and $D$, however, is arbitrary insofar as it depends on the choice of $e$, which itself always belongs to $B$. Thus, if a number from $D$ is chosen instead, the classes $B,D$ are interchanged.

\Article{6}
Since $-1$ is a quadratic residue of $p$, let us put $-1\equiv ff\; (\mathrm{mod.}\,p)$. The four roots of the congruence $x^4\equiv1$ are then $1,f,-1,-f$. Thus if\space%
\SourceLeaf{81}%
a is a biquadratic residue of $p$, say $\equiv\alpha^4$, the four roots of the congruence $x^4\equiv a$ are $\alpha,f\alpha,-\alpha,-f\alpha$, which are easily seen to be pairwise incongruent. Hence, if we collect the least positive residues of the fourth powers $1,16,81,256\dotsfour(p-1)^4$, they are always equal in groups of four, giving $\tfrac14(p-1)$ distinct biquadratic residues that form the collection $A$. If we collect the least residues of the fourth powers only as far as $(\tfrac12p-\tfrac12)^4$, each occurs twice.

\Article{7}
The product of two biquadratic residues is clearly a biquadratic residue: multiplying two numbers from class $A$ always gives a product whose least positive residue belongs to the same class. Similarly, the product of a number from $B$ and one from $D$, or of a number from $C$ and one from $C$, has its least residue in $A$.

In $B$ lie the residues of the products $A.B$ and $C.D$; in $C$ those of $A.C$, $B.B$ and $D.D$; and finally in $D$ those of $A.D$ and $B.C$.

The proofs are so evident that it suffices to indicate one. Let, for example, $c$ and $d$ be numbers from $C$ and $D$, and put $c\equiv eea$, $d\equiv e^3a'$, with $a,a'$ numbers from $A$. Then $e^4aa'$ is a biquadratic residue, so its least residue belongs to $A$. For the product we have $cd$\relax$\equiv e.e^4aa'$, so its least residue belongs to $B$.

We can now also easily decide which class contains a product of several factors. Assign to the classes $A,B,C,D$ the characters $0,1,2,3$, respectively. The character of a product is then either the sum of the characters of its individual factors or the least residue of that sum modulo $4$.

\Article{8}
It seemed worthwhile to develop these elementary propositions without using the theory of power residues. With its help, everything can be proved much more easily still.

\Needspace{16\baselineskip}Let $g$ be a primitive root for the modulus $p$: a number such that in the sequence of powers $g,gg,g^3\dotsfour$, none before $g^{p-1}$ is congruent to unity modulo $p$. Then the least positive residues of the numbers $1,g,gg,g^3\dotsfour g^{p-2}$ agree, apart from their order, with $1,2,3\dotsfive p-1$, and are distributed among four classes as follows:

\SourceLeaf{82}%
\[
\begin{array}{c|l}
\text{to}&\text{least residues of the numbers}\\\hline
A&1,\ g^4,g^8,g^{12}\dotsfive g^{p-5}\\
B&g,\ g^5,g^9,g^{13}\dotsfive g^{p-4}\\
C&gg,g^6,g^{10},g^{14}\dotsfive g^{p-3}\\
D&g^3,g^7,g^{11},g^{15}\dotsfive g^{p-2}
\end{array}
\]

All the preceding propositions follow directly from this.

Just as the numbers $1,2,3\dotsfour p-1$ have here been divided among four classes, whose collections we denote by $A,B,C,D$, we may assign \emph{any} integer not divisible by $p$ to one of these classes according to its least residue modulo $p$.

\Article{9}
We denote by $f$ the least residue of the power $g^{\frac14(p-1)}$ modulo $p$. Since $ff\equiv g^{\frac12(p-1)}\equiv-1$ (\emph{Disquis. Arithm.}, Article 62), the symbol $f$ clearly has the same meaning here as in Article 6. Thus the power $g^{\frac14\lambda(p-1)}$, with $\lambda$ a positive integer, is congruent modulo $p$ to $1,f,-1,-f$, according as $\lambda$ has the form $4m,4m+1,4m+2,4m+3$, respectively, or according as the least residue of $g^\lambda$ belongs to $A,B,C,D$, respectively. This gives us a very simple criterion for deciding which class contains a given number $h$ not divisible by $p$: $h$ belongs to $A,B,C$ or $D$ according as the power $h^{\frac14(p-1)}$ is congruent modulo $p$ to $1,f,-1$ or $-f$.

As a corollary, $-1$ always belongs to class $A$ when $p$ has the form $8n+1$, and to class $C$ when $p$ has the form $8n+5$. A proof of this theorem independent of the theory of power residues can easily be obtained from what we established in Article 115, III of the \emph{Disquisitiones Arithmeticae}.

\Article{10}
Since \emph{all} primitive roots for the modulus $p$ arise as residues of the powers $g^\lambda$, on taking for $\lambda$ all numbers coprime to $p-1$, it is easy to see that these roots are equally divided between $B$ and $D$, with the base $g$ always in $B$. If we replace $g$ as base by another primitive root from $B$, the classification remains the same; if we use a primitive root from $D$ as base, the classes $B$ and $D$ are interchanged.

\SourceLeaf{83}%
If the classification is based on the criterion in the preceding article, the distinction between $B$ and $D$ depends on which root of the congruence $xx\equiv-1\; (\mathrm{mod.}\,p)$ we choose as the characteristic number $f$.

\Article{11}
To make it easier to illustrate by examples the subtler investigations that we are about to undertake, we give here the construction of the classes for all moduli below $100$. For each modulus we have chosen the least primitive root.

\[
\begin{gathered}
p=5\\
g=2,\ f=2\\[3pt]
\begin{array}{c|l}
A&1\\
B&2\\
C&4\\
D&3
\end{array}
\end{gathered}
\]

\[
\begin{gathered}
p=13\\
g=2,\ f=8\\[3pt]
\begin{array}{c|l}
A&1,3,9\\
B&2,5,6\\
C&4,10,12\\
D&7,8,11
\end{array}
\end{gathered}
\]

\[
\begin{gathered}
p=17\\
g=3,\ f=13\\[3pt]
\begin{array}{c|l}
A&1,4,13,16\\
B&3,5,12,14\\
C&2,8,9,15\\
D&6,7,10,11
\end{array}
\end{gathered}
\]

\[
\begin{gathered}
p=29\\
g=2,\ f=12\\[3pt]
\begin{array}{c|l}
A&1,7,16,20,23,24,25\\
B&2,3,11,14,17,19,21\\
C&4,5,6,9,13,22,28\\
D&8,10,12,15,18,26,27
\end{array}
\end{gathered}
\]

\SourceLeaf{84}%
\[
\begin{gathered}
p=37\\
g=2,\ f=31\\[3pt]
\begin{array}{c|l}
A&1,7,9,10,12,16,26,33,34\\
B&2,14,15,18,20,24,29,31,32\\
C&3,4,11,21,25,27,28,30,36\\
D&5,6,8,13,17,19,22,23,35
\end{array}
\end{gathered}
\]

\[
\begin{gathered}
p=41\\
g=6,\ f=32\\[3pt]
\begin{array}{c|l}
A&1,4,10,16,18,23,25,31,37,40\\
B&6,14,15,17,19,22,24,26,27,35\\
C&2,5,8,9,20,21,32,33,36,39\\
D&3,7,11,12,13,28,29,30,34,38
\end{array}
\end{gathered}
\]

\[
\begin{gathered}
p=53\\
g=2,\ f=30\\[3pt]
\begin{array}{c|l}
A&1,10,13,15,16,24,28,36,42,44,46,47,49\\
B&2,3,19,20,26,30,31,32,35,39,41,45,48\\
C&4,6,7,9,11,17,25,29,37,38,40,43,52\\
D&5,8,12,14,18,21,22,23,27,33,34,50,51
\end{array}
\end{gathered}
\]

\[
\begin{gathered}
p=61\\
g=2,\ f=11\\[3pt]
\begin{array}{c|l}
A&1,9,12,13,15,16,20,22,25,34,42,47,56,57,58\\
B&2,7,18,23,24,26,30,32,33,40,44,50,51,53,55\\
C&3,4,5,14,19,27,36,39,41,45,46,48,49,52,60\\
D&6,8,10,11,17,21,28,29,31,35,37,38,43,54,59
\end{array}
\end{gathered}
\]

\[
\begin{gathered}
p=73\\
g=5,\ f=27\\[3pt]
\begin{array}{c|l}
A&1,2,4,8,9,16,18,32,36,37,41,55,57,64,65,69,71,72\\
B&5,7,10,14,17,20,28,33,34,39,40,45,53,56,59,63,66,68\\
C&3,6,12,19,23,24,25,27,35,38,46,48,49,50,54,61,67,70\\
D&11,13,15,21,22,26,29,30,31,42,43,44,47,51,52,58,60,62
\end{array}
\end{gathered}
\]

\SourceLeaf{85}%
\[
\begin{gathered}
p=89\\
g=3,\ f=34\\[3pt]
\begin{array}{c|l}
A&1,2,4,8,11,16,22,25,32,39,44\mathpunct{.}\,45,\\
&50,57,64,67,73,78,81,85,87,88\\
B&3,6,7,12,14,23,24,28,33,41,43,46,\\
&48,56,61,65,66,75,77,82,83,86\\
C&5,9,10,17,18,20,21,34,36,40,42,47,\\
&49,53,55,68,69,71,72,79,80,84\\
D&13,15\mathpunct{.}\,19,26,27,29\mathpunct{.}\,30,31,35,37,38\mathpunct{.}\,51,\\
&52,54,58,59,60,62,63,70,74,76
\end{array}
\end{gathered}
\]

\[
\begin{gathered}
p=97\\
g=5,\ f=22\\[3pt]
\begin{array}{c|l}
A&1,4,6,9,16,22,24,33,35,36,43,47,\\
&50,54,61,62,64,73,75,81,88,91,93,96\\
B&5,13,14,17,19,20,21,23,29,30,41,45,\\
&52,56,67,68,74,76,77,78,80,83,84,92\\
C&2,3,8,11,12,18,25,27,31,32,44,48,\\
&49,53,65\mathpunct{.}\,66,70,72,79,85,86,89,94,95\\
D&7,10,15,26,28,34,37,38,39\mathpunct{.}\,40,42,46,\\
&51,55,57,58,59,60,63,69,71,82,87,90
\end{array}
\end{gathered}
\]

\Article{12}
Since $2$ is a quadratic residue of every prime of the form $8n+1$, but a nonresidue of every prime of the form $8n+5$, for prime moduli of the first form $2$ belongs to class $A$ or $C$, and for moduli of the second form to $B$ or $D$. The distinction between $B$ and $D$ is not intrinsic, depending only on the choice of $f$, so we shall set moduli of the form $8n+5$ aside for a while. An \emph{inductive investigation} of moduli of the form $8n+1$ gives $2$ in $A$ for $p=73,89,113,233,257,281,337,353$, etc.; in contrast, $2$ belongs to $C$ for $p=17,41,97,137,193,241,313,401,409,433,449,457$, etc.

Since, moreover, for a prime modulus of the form $8n+1$, $-1$ is a biquadratic residue, $-2$ clearly always belongs to the same class as $+2$.
